\section{Zephyr、Tulu 2、SteerLM 与 Starling：没有单一赢家}

上一章给出公式，本章看模型证据。Zephyr 证明 DPO 能做出出圈开放模型，Tulu 2 把 DPO 扩展到 70B；SteerLM 与 Starling 又证明 PPO/RL 路线仍可能更强。历史结论不是“DPO 取代 PPO”，而是 preference optimization 开始产生可比较的公开 releases。

\subsection{Zephyr 与 Tulu 2：DPO 从论文变成 release}

本节先看两次关键落地。Zephyr 结合 UltraChat 与 UltraFeedback，给 DPO 一个清晰、可复现 recipe；Tulu 2 使用 AI2 数据与训练栈，把方法扩展到 70B。模型发布让社区能检验方法是否跨规模、跨数据保持价值。

\lecturefigure{slide-067.jpg}{Zephyr $\beta$：第一批因 DPO 获得广泛关注的开放模型}{official deck, p.~67}

\lecturefigure{slide-068.jpg}{Tulu 2：把 DPO 扩展到 70B 规模}{official deck, p.~68}

\begin{knowledgebox}{读图：recipe 的可迁移部分}
比较 release 时要拆成 base model、IFT data、preference data、reference policy、$\beta$、sequence length 和 eval。若 Zephyr 与 Tulu 同时换了多个组件，不能把全部提升归因于 DPO；它们的贡献是给出完整 artifact，让后续团队能逐项消融。
\end{knowledgebox}

\teachervoice{00:46:33--00:48:30，Nathan 把 Zephyr 称为第一个让 DPO 真正“make a splash”的模型，把 Tulu 2 视为规模化里程碑。方法的社会可信度来自模型 release，而不只是论文推导。}

\subsection{SteerLM 与 Starling：PPO/RL 仍然竞争}

DPO 简单并不意味着 on-policy RL 失效。SteerLM 用属性条件化和偏好信号控制输出，Starling 基于 Arena 风格偏好数据运行 PPO，展示更复杂的 pipeline 仍能获得优势。课程特意把这些模型放在 DPO 成功之后，阻止“新方法统一一切”的叙事。

\lecturefigure{slide-069.jpg}{SteerLM 与 Starling：PPO/RL 方法仍能超过 DPO 模型}{official deck, p.~69}

\begin{warningbox}{算法榜单常把数据优势伪装成优化器优势}
若 PPO 模型使用更新、更大或更真实的 preference data，而 DPO baseline 使用旧数据，结果不能归因于 optimizer。可信比较需要共享 base/IFT checkpoint、共享 pairs 或 reward model、相似计算预算和同一 eval portfolio。
\end{warningbox}

\teachervoice{00:48:45--00:49:35，讲者用 SteerLM 和 Starling 说明“仍有很多模型显示 PPO/RL 优于 DPO”。他的结论是混合证据，而不是为某一阵营站队。}

\subsection{本章小结}

Zephyr 与 Tulu 2 证明 DPO 的实用性和规模化，SteerLM/Starling 保留了 PPO 的竞争力。开放 alignment 的合理实验单位不是算法名字，而是完整 recipe。下一章进入 2024 年生态，讨论更多参与者、Llama 3、开放/闭源差距和真正的瓶颈——高质量 preference data。

